\begin{center}
{\zihao{2}\bfseries 考虑变物性与径向收缩的药材\\[5pt]热风干燥建模及数值分析}
\end{center}
\vspace{3pt}
\begin{center}\bfseries 摘\quad 要\end{center}
热风干燥中药材内部温湿分布不均、传递参数随状态变化及失水收缩共同影响干燥终点判断。本文根据给定环境数据与经验物性，建立圆柱径向有效传热传质模型，以守恒型有限体积法和自适应隐式积分求解四个问题。

\textbf{针对问题一，建立常热物性下的非线性扩散模型。}由圆柱薄壳的能量和水质量收支推导控制方程，以中心对称和表面对流交换闭合初边值问题；对环境数据分段线性插值，采用表面加密网格处理近表面水分梯度。1800 s时，中心、表面温度分别为\textbf{33.5753、36.7856 $^\circ$C}，干基含水率分别为\textbf{2.5500、1.5102 kg/kg}，并得到每秒、每0.1 cm的分布。

\textbf{针对问题二，建立由含水率和温度共同驱动的变物性耦合模型。}将经验密度、比热容、导热系数和扩散系数嵌入局部传递方程，通过共享界面通量和稀疏雅可比矩阵实现温湿状态的同步隐式更新。3 h时，中心、表面温度分别为\textbf{49.8495、49.9664 $^\circ$C}，含水率分别为\textbf{1.7662、1.0081 kg/kg}，说明热平衡明显早于水分达标。

\textbf{针对问题三，以全域最大含水率构造干燥终点事件。}在问题二的模型上延长积分，将4 h之后环境取为3--4 h观测均值，并用连续数值解定位最大含水率达到0.15 kg/kg的时刻。得到固定半径下干燥临界时长为\textbf{57.4740 h}，临界时中心为最后达标区域，表面含水率为0.0526 kg/kg。

\textbf{针对问题四，建立随材料运动的收缩域模型。}由干物质骨架与水分连续性方程出发，引入均匀径向收缩和坐标变换$x=r/R(t)$，将移动边界问题转化到固定区间，再代入题给收缩数据及新物性。计算得到干燥临界时长为\textbf{51.0920 h}，此时半径为1.2000 cm。该值比问题三减少约6.3820 h，是物性变化与几何变化共同作用的结果。

采用贝塞尔级数独立解、空间加密、时间容差收紧及整体守恒检查核验既有计算：问题一温度与级数解的最大差约$2.80\times10^{-7}\,^\circ$C；末两级网格的问题三、四时长差分别为0.0243 s、0.0095 s。上述结果适用于本文明确的有效湿度边界、无显式潜热及后期环境延长假设。

\vspace{5pt}
\noindent\textbf{关键词：}药材干燥；热质耦合；有限体积法；移动边界；自适应BDF

\nextpage
\section{问题背景与重述}
\subsection{问题背景}
中药材采收后含有较多水分，干燥过程既关系到储存稳定性，也影响后续加工质量。热风干燥通过提高环境温度、维持空气与材料之间的传热传质驱动力，使内部水分迁移到表面并进入周围空气。其工艺通常经历预热平衡与恒温干燥两个阶段：前者以内外温差引起的温度调整为主，后者即使外部温度基本稳定，内部水分仍可能持续变化。

对于具有一定厚度的药材，表面直接与空气交换，中心则通过内部传递间接受到影响。因而单一平均温度或平均含水率不能充分描述产品各处是否达到要求。另一方面，扩散能力和热物性会随温湿状态变化，失水还可能导致几何尺寸收缩；若始终采用固定参数和固定尺寸，得到的局部分布及终点时间可能出现系统偏差。干燥理论通常通过守恒方程与材料传递关系构建可计算模型\cite{mujumdar,whitaker,defraeye}，其优势在于能够把局部机制、外部条件与全过程结果联系起来。

本题给出的圆柱形药材长度为25 cm、初始半径为2 cm，初始温度为28 $^\circ$C、初始干基含水率为2.55 kg/kg。附件提供早期烘房温湿数据与随时间变化的半径，题目附录提供各阶段使用的物性参数。基于这些数据，本文研究径向温湿演化和全域干燥终点，并对已完成的数值结果作一致性分析。

\subsection{问题重述}
\textbf{问题一：预热阶段的分布预测。}采用题目附录2的常热物性及含水率相关扩散系数，预测前1800 s的温度、干基含水率分布；在指定七个时刻与五个半径位置报告结果，并提供每隔1 s、0.1 cm的完整输出。

\textbf{问题二：全过程的变物性描述。}统一采用题目附录3的经验关系，建立温度与含水率相互作用的动态模型；报告前3 h中每隔0.5 h的五个径向位置结果，并形成每秒的径向分布记录。

\textbf{问题三：固定尺寸下的干燥时长。}沿用问题二模型，以材料各处含水率均低于0.15 kg/kg为目标，确定干燥所需时间；按6 h间隔列出指定位置含水率，并保存每60 s、每0.1 cm的明细。

\textbf{问题四：考虑收缩后的干燥过程。}结合附件2中的半径变化及题目附录4的新物性，建立收缩材料域上的传递模型，确定干燥时长，并报告实际空间位置与当前表面的水分分布。

\note{文中$C$均指药材干基含水率。空气数据如何作为材料表面的有效平衡边界，将在模型假设及问题一中明确；题目提供的数值本身不构成完整的吸附等温关系。}

\nextpage
\section{问题分析}
\textbf{问题一分析。}研究对象是具有径向空间差异的圆柱体，需要同时获得$T(r,t)$和$C(r,t)$，不能只建立整体平均量的常微分方程。对任意圆柱薄壳分别写能量、水质量守恒，再用傅里叶定律和有效菲克关系表达界面通量，即可得到径向扩散型偏微分方程。中心由对称性约束，表面由对流交换约束，环境附件只进入表面条件。第一问热物性固定，温度方程为线性方程；扩散系数随$C$变化，水分方程仍为非线性方程，适合用守恒离散统一处理。

\textbf{问题二分析。}密度、比热容、导热系数随含水率变化，水分扩散又受到绝对温度影响，因而水分和温度之间形成双向反馈。若先独立求完温度，再把温度代入水分方程，只能得到弱耦合近似。本文将两类状态交错组成一个未知向量，在同一隐式时间步内更新物性和通量。为与题给经验式及既有算法一致，密度用于有效体积热容量，水分以均匀参考干质量为基准描述；不再同时施加可能不相容的精确湿质量密度约束。

\textbf{问题三分析。}“各处低于阈值”是全域条件，不等价于平均含水率达标或表面达标。应对全部径向节点取最大值，构造阈值事件，并在连续数值轨迹上定位穿越时刻。数据仅覆盖4 h，长达两三天的计算还必须给出后期环境延长规则。本文采用稳定阶段均值，并利用已有末值保持方案检验这一选择的影响；同时区分达到0.15的临界时间与严格低于0.15的后续时刻。

\textbf{问题四分析。}半径变化不仅改变扩散距离，也改变材料点与固定空间点的对应关系。直接把定域方程中的$R$替换为$R(t)$而不说明材料运动，难以保证守恒解释一致。本文先建立骨架与水分连续性方程，假设均匀径向收缩，利用$x=r/R(t)$转到随骨架运动的坐标；随后推导扩散项及交换边界的半径因子，计算结束时再还原到实际距离。对已经处于材料域外的位置不赋予虚构含水率。

\begin{figure}[H]\centering
\begin{tikzpicture}[node distance=8mm and 11mm]
\node[box] (a) {附件与物性\\环境、半径、经验式};
\node[box,right=of a] (b) {局部守恒建模\\通量、初值、边界};
\node[box,right=of b] (c) {有限体积离散\\稀疏隐式积分};
\node[box,below=of c] (d) {分布与阈值事件\\四问结果输出};
\node[box,left=of d] (e) {独立解与收敛\\守恒及敏感性};
\node[box,left=of e] (f) {条件化结论\\工艺解释与改进};
\draw[arr](a)--(b);\draw[arr](b)--(c);\draw[arr](c)--(d);
\draw[arr](d)--(e);\draw[arr](e)--(f);
\end{tikzpicture}
\caption{建模、求解与验证的总体技术路线}\label{fig:roadmap}
\end{figure}

\nextpage
\section{模型假设}
为使给定数据能够闭合为可计算问题，并保持论文与既有算法一致，采用以下假设。

\begin{enumerate}[label=\textbf{假设\arabic*：},leftmargin=4.7em]
\item \textbf{圆柱径向主导。}药材初始形状近似为圆柱，材料及外部交换在周向上均匀，忽略轴向梯度与端面传递，状态仅依赖半径和时间。端面总面积与侧面积之比为$R/L=0.08$，这一几何关系支持径向近似作为初步模型，但不意味着端部局部效应严格为零。
\item \textbf{连续有效介质。}用宏观温度和干基含水率概括孔隙与固体的平均状态；局部热流服从傅里叶关系，水分通量以含水率梯度驱动。经验系数直接采用题给关系，不另外识别微观参数。
\item \textbf{显热近似。}温度方程以$\rho c_p$表示有效体积热容量，忽略显式蒸发潜热、迁移水分携带的焓、辐射以及反应热；表面不设独立热储存。故本文模拟的是题给参数下的有效热质传递过程，未建立完整多相能量平衡。
\item \textbf{有效湿度边界。}附件1的空气水分数值直接作为与给定传质系数配套的有效平衡干基值$C_e(t)$。该关系是闭合假设，不将空气含湿比与固体干基含水率视为一般意义上可直接等同的物理量。
\item \textbf{密度作用分开约定。}题给$\rho(C)$用于热方程；水分按均匀参考干密度$\rho_d$守恒。前三问无骨架运动，$\rho_d$在时间和空间上固定；第四问允许其随整体体积变化。本文不额外要求$\rho(C)=\rho_d(1+C)$处处精确成立。
\item \textbf{环境与交换条件。}4 h内的环境数据线性插值，4 h后以3--4 h数据均值保持；各问均采用$h=25\ \mathrm{W/(m^2K)}$和$h_m=8\times10^{-7}$ m/s。附件没有单独给出的换热、传质系数变化不另行假定。
\item \textbf{均匀径向收缩。}第四问中长度固定，外半径由附件2线性插值；内部骨架速度取$v=r\dot R/R$，不考虑开裂、分层和不均匀变形。固定$x=r/R(t)$对应同一材料位置。
\end{enumerate}

上述假设使四问使用一致的物理和数值框架。前三问在固定域上比较不同物性与终点条件，第四问再增加测得几何变化。对于题目没有给出的本构关系，本文明确保留为模型条件；对于已经给出的观测输入，则尽可能直接使用，避免通过额外拟合引入没有数据支持的自由参数。

\note{数值收敛检查回答“所列方程是否被准确求解”，不直接证明全部物理假设均成立。后文分别报告离散误差和环境延长方式的敏感性。}

\nextpage
\section{符号说明}
\begin{table}[H]\centering
\caption{主要符号及其单位}\label{tab:symbols}
\small
\begin{tabularx}{\linewidth}{>{\centering\arraybackslash}p{24mm}X>{\centering\arraybackslash}p{34mm}}
\toprule
符号 & 说明 & 单位\\\midrule
$r,\ R(t),\ L$ & 径向坐标、外半径、圆柱长度 & m\\
$t,\ t_*$ & 时间、干燥临界时间 & s\\
$T,\ T_a,\ T_s$ & 药材、空气及表面温度 & $^\circ$C\\
$T_K$ & 计算经验扩散系数使用的绝对温度 & K\\
$C,\ C_s,\ C_e$ & 内部、表面及有效平衡干基含水率 & kg/kg\\
$m_d,\ m_w$ & 干物质量、水质量 & kg\\
$\rho,\ \rho_d$ & 有效热密度、单位体积参考干物质量 & kg/m$^3$\\
$c_p$ & 有效比热容 & J/(kg$\cdot$K)\\
$k$ & 导热系数 & W/(m$\cdot$K)\\
$D$ & 有效水分扩散系数 & m$^2$/s\\
$h$ & 表面对流换热系数 & W/(m$^2\cdot$K)\\
$h_m$ & 表面有效传质系数 & m/s\\
$q,\ j$ & 向外热流密度、相对骨架水质量通量 & W/m$^2$；kg/(m$^2\cdot$s)\\
$v$ & 径向骨架速度 & m/s\\
$x,\ x_i$ & 归一化坐标及其节点位置 & 1\\
$N,\ \beta$ & 径向网格区间数、加密参数 & 1\\
$\omega_i$ & 归一化控制体积的径向权重 & 1\\
$H_{i+1/2}$ & 归一化热通量，定义见式\eqref{eq:fluxes} & W/m$^3$\\
$W_{i+1/2}$ & 归一化含水率通量 & s$^{-1}$\\
$\alpha,\ Bi$ & 热扩散率、基于半径的热毕渥数 & m$^2$/s；1\\
$M(t),\ \overline C(t)$ & 全域最大、干质量加权平均含水率 & kg/kg\\
$\bm y,\ \bm F$ & 离散状态向量、其时间变化率 & 分量单位各异\\
\bottomrule
\end{tabularx}
\end{table}

\textbf{单位与基准约定。}内部计算统一采用米和秒，输出时再转为厘米、小时。干基含水率定义为$C=m_w/m_d$，数值可大于1；例如$C=2.55$对应湿基质量分数$2.55/(1+2.55)\approx71.83\%$，不能与0.15的干基阈值交叉比较。

\textbf{温度与导数约定。}温差使用摄氏度与开尔文度数差完全等价，但指数经验式中的温度必须取$T_K=T+273.15$。下标$r,t,x$表示相应偏导，点号表示时间导数；表面下标$s$表示在$r=R(t)$处取值。公式中的通量均以从中心向外为正。
